\section{经济性与模型组合策略}

\begin{importantbox}{强弱模型组合原则}
\begin{itemize}
    \item \textbf{Planning 和 RePlan}：使用强模型（如推理模型），因为规划决定了整个 Workflow 的上限
    \item \textbf{Execute}：使用弱模型（如 GPT-4.1 Nano/Mini），因为已经告诉模型该怎么做，只需一步步执行
\end{itemize}
这与 OpenAI 对推理模型和非推理模型的定位完全一致。
\end{importantbox}

如果 Planning 环节特别重要，还可以引入 \textbf{Generator-Critique} 模式来形成更稳固可靠的 Plan——不要预期模型一次性就能提出一个很好的 Plan。

\subsection{本章小结}

合理的模型组合策略可以在成本和效果之间取得平衡：贵的模型做规划，便宜的模型做执行。

